
\medskip

\textbf{Question 1}

Quel est le tiers de 18 ?

\medskip

\textbf{Question 2}

Un film dure $240$~min. Quelle est sa durée en heures ?

\medskip

\textbf{Question 3}

Les notes obtenues par un élève sont : 8~;~12~;~6~;~19~;~15

Que vaut la médiane de cette série de notes ?

\medskip

\textbf{Question 4}

\begin{center}
\psset{unit=1cm}
\begin{pspicture}(0,-1)(9.,1)
\psgrid[gridlabels=0pt,subgriddiv=4,gridwidth=0.25pt,subgridwidth=0.25pt]
\psline(9.2,0)
\qdisk(2,0){2pt}\qdisk(6,0){2pt}\qdisk(9,0){2pt}
\uput[d](2,0){0}\uput[d](6,0){1}\uput[u](9,0){E}
\end{pspicture}
\end{center}

Sur cette droite graduée, l'abscisse du point E est
\begin{center}
\begin{tabularx}{\linewidth}{*{4}{X}}
\textbf{A.~} $\dfrac54$ &\textbf{B.~} $\dfrac32$ &\textbf{C.~} $\dfrac74$ &\textbf{D.~} $\dfrac52$\\ 
\end{tabularx}
\end{center}

\medskip

\textbf{Question 5}

\begin{minipage}{0.48\linewidth}
Dans le triangle ABC, rectangle en B,
on sait que $\widehat{\text{A}} = 35\degres$.

Calculer $\widehat{\text{C}}$.
\end{minipage}\hfill
\begin{minipage}{0.48\linewidth}
\psset{unit=1cm}
\begin{pspicture}(5,3)
\pspolygon(0.2,0)(4.6,0)(4.6,2.6)%ABC
\psframe(4.6,0)(4.4,0.2)
\uput[ul](0.2,0){A} \uput[ur](4.6,0){B} \uput[ur](4.6,2.6){C} \uput[ur](0.9,0.1){$35\degres$} \uput[dl](4.6,2.6){?}
\psarc(0.2,0){0.5}{0}{30.58}\psarc(4.6,2.6){0.5}{211.58}{270}
\end{pspicture}
\end{minipage}

\medskip

\textbf{Question 6}

\begin{minipage}{0.6\linewidth}
Dans le triangle ABC, rectangle en A, quel calcul doit-on effectuer pour déterminer le cosinus de l'angle 
$\widehat{\text{ABC}}$ ?
\end{minipage}\hfill
\begin{minipage}{0.38\linewidth}
\psset{unit=1cm}
\begin{pspicture}(5.2,2.3)
\pspolygon(0.2,0)(5.2,0)(0.2,2.1)%ABC
\psframe(0.2,0)(0.4,0.2)
\uput[ul](0.2,0){A} \uput[ur](5.2,0){B} \uput[ul](0.2,2.1){C} 
\end{pspicture}
\end{minipage}

\medskip

\textbf{Question 7}

\begin{minipage}{0.45\linewidth}
Sur la figure ci-contre, dans le triangle ADE les droites (DE) et (CB) sont parallèles.

Déterminer la longueur AD.
\end{minipage}\hfill
\begin{minipage}{0.52\linewidth}
\psset{unit=0.9cm}
\begin{pspicture}(8,3.5)
\pspolygon(0.2,0.2)(7.7,0.2)(6.2,3.2)%AED
\psline(2.4,0.2)(2,1.1)%BC
\uput[ul](0.2,0.2){A} \uput[ur](7.7,0.2){E} \uput[ur](6.2,3.2){D} 
\uput[d](2.4,0.2){B}\uput[ul](2,1.1){C}
\uput[ur](2.2,0.6){2 cm} \uput[ul](1.2,0.7){4 cm} \uput[ur](6.95,1.7){7 cm}
\end{pspicture}
\end{minipage}

\medskip

\textbf{Question 8}

Dans un collège, 25\,\% des 300 élèves participent à une olympiade de mathématiques.

Combien d'élèves ne participent pas a cette olympiade ?

\medskip


\begin{minipage}{0.6\linewidth}
\textbf{Question 9}

Une élève souhaite réaliser un programme avec un logiciel de
programmation pour dessiner un carré.

Par quelles valeurs doit-on compléter les lignes 3 et 5 
pour obtenir un carré ?
\end{minipage}\hfill
\begin{minipage}{0.34\linewidth}
\begin{scratch}[scale=0.8,num blocks, num start=1,baseline=2]]
\initmoreblocks{definir carré}
\blockpen{stylo en position d'écriture}
\blockrepeat{répéter \ovalnum{\phantom{aa}} fois}
{
\blockmove{avancer de \ovalvariable{50} pas}
\blockmove{tourner \turnleft{} de \ovalnum{\phantom{aa}} degrés}
}
\end{scratch}
\end{minipage}


%%%%%% Partie avec calculatrice
% \begin{center}
% \textbf{Partie 2 — Raisonnement et résolution de problèmes — 14 points — 1 h 40}
% \end{center}

% \medskip

% \begin{tabular}{|p{14cm}|}\hline
% \textbf{Dans cette partie, toutes les réponses doivent être justifiées, sauf si une indication contraire est donnée.}\\
% \textbf{La clarté et la précision des raisonnements ainsi que la rédaction sont évaluées sur 2 points.}\\
% \textbf{Pour chaque question, si le travail n'est pas terminé, laisser tout de même une trace de la recherche ; les essais et les démarches engagées, même non aboutis, seront pris en compte dans la notation.}\\ \hline
% \end{tabular}

% \medskip

